% Colonial English Phonology as a Filter — Thematic Synthesis
% !TEX program = xelatex
% !TEX encoding = UTF-8 Unicode

\section{Colonial English Phonology as a Filter}
\label{sec:colonial-phonology}

Each recorder's English orthography is a lens through which Algonquian
phonemes were filtered. Understanding the recorder's native phonological
system is essential to recovering the target language's sound system. The
colonial recorders did not write what they heard --- they wrote what their
English phonological categories allowed them to perceive. This section
establishes the Early Modern English baseline for each recorder and
documents how their native systems systematically distorted their
transcriptions.

\subsection{Early Modern English Phonology Baseline (c.\ 1600--1620)}

The Great Vowel Shift (GVS) was complete by approximately 1600--1620 for
most English dialects. This means the colonial recorders' vowel system was
significantly different from both Middle English and Modern English.
Table~\ref{tab:eme-vowels} shows the vowel values that the earliest
recorders (Winslow, Wood, Williams) would have used.

\begin{table}[h]
\centering
\caption{Early Modern English vowel values (c.\ 1600--1620)}
\label{tab:eme-vowels}
\begin{tabular}{llll}
\toprule
\textbf{Spelling} & \textbf{ME value} & \textbf{c.\ 1600 value} & \textbf{Modern value} \\
\midrule
\textit{a} (short) & /a/ & /a/ & /æ/ \\
\textit{e} (short) & /ɛ/ & /ɛ/ & /ɛ/ \\
\textit{i} (short) & /ɪ/ & /ɪ/ & /ɪ/ \\
\textit{o} (short) & /ɔ/ & /ɔ/ & /ɒ/~/ɑ/ \\
\textit{u} (short) & /ʊ/ & /ʌ/ & /ʌ/ \\
\textit{a} (long) & /aː/ & /ɛː/~/eː/ & /eɪ/ \\
\textit{ee} (long) & /eː/ & /iː/ & /iː/ \\
\textit{ie} (long) & /ɛː/ & /iː/ & /iː/ \\
\textit{oo} (long) & /oː/ & /uː/ & /uː/ \\
\textit{ou} (long) & /uː/ & /aʊ/ & /aʊ/ \\
\textit{ea} (digraph) & /ɛː/ & /ɛː/~/eː/ & /iː/ \\
\textit{ai/ay} & /aɪ/ & /ɛː/~/eː/ & /eɪ/ \\
\bottomrule
\end{tabular}
\end{table}

The critical point for Algonquian transcription is that the \textit{ea}
digraph in 1620--1643 represented /ɛː/ or /eː/, NOT the modern /iː/. When
Williams writes \textit{wean} for a Narragansett word, the \textit{ea}
represents a mid front vowel, not a high front vowel. Similarly,
\textit{oo} in 1620 represented /oː/ (still raising toward /uː/), not the
modern /uː/. These differences matter for mapping colonial spellings to
Algonquian phonemes.

\subsection{Per-Recorder Linguistic Background}

Each recorder brought a different linguistic background to the task of
transcribing an unfamiliar language. Table~\ref{tab:recorder-background}
summarizes their training and its impact.

\begin{table}[h]
\centering
\caption{Recorder linguistic background and transcription impact}
\label{tab:recorder-background}
\begin{tabular}{llll}
\toprule
\textbf{Recorder} & \textbf{Training} & \textbf{Key bias} & \textbf{Impact} \\
\midrule
Winslow (1624) & None (colonist) & English-hearer perception & Raw phonetic impression \\
Wood (1634) & None (colonist) & English-hearer perception & b/p alternation, unstable vowels \\
Williams (1643) & Theologian, self-taught linguist & English categories + stress awareness & Diacritics for stress, no /h/ \\
Eliot (1666) & Missionary, Hebrew/Greek training & Systematic but English-filtered & Aspiration via voiced/voiceless \\
Edwards (1787) & Theologian, fluent L2 speaker & Hebrew/Greek + childhood fluency & Most reliable /h/ detection \\
Prince (1904) & Neogrammarian linguist & German orthographic conventions & \textit{tsh} for /tʃ/, \textit{Dehnungs-h} \\
Speck (1904) & Boasian anthropologist & Ethnographic, minimal linguistic training & Two-stage recording chain \\
Fielding (2013) & Modern linguist (MIT) & Phonemic orthography & 1:1 phoneme-to-grapheme \\
\bottomrule
\end{tabular}
\end{table}

\subsubsection{Winslow (1624) --- The Honest Amateur}

Edward Winslow was unusually honest about his linguistic limitations,
describing the Massachusett language as ``very copious, large, and
difficult.'' His transcriptions are raw phonetic impressions --- he was not
normalizing or regularizing what he heard, unlike Eliot (who developed a
systematic orthography) or Williams (who, though linguistically gifted,
worked through an English lens). Winslow's data passed through two
intermediaries: Tisquantum (Squanto), a Patuxet Wampanoag who learned
English in London and may have experienced language attrition, and
Hobbamock, a Pokanoket \textit{pniese} (military champion) with limited
English. This intermediary filter means every transcription is at least two
steps removed from the speaker.

\subsubsection{Wood (1634) --- The Pre-Orthographic Recorder}

William Wood's orthography is pre-Eliot --- the Massachusett writing system
was still decades away. His transcriptions reflect raw English-hearer
perception without regularization. The most diagnostic feature is his b/p
alternation: Wood writes \textit{abbona} (five) and \textit{abamocho} (devil)
with \textit{b} where the Eliot/Trumbull orthography uses \textit{p}. This
shows Wood perceived the unaspirated Wampanoag /p/ as closer to English /b/
--- direct evidence of the voicing neutralization that characterizes SNEA
stop systems. Wood's triple spelling for a single word
(\textit{abbomocho}/\textit{abomocho}/\textit{abamacho}) is the strongest
evidence in the entire SNEA corpus for colonial vowel instability.

\subsubsection{Williams (1643) --- The Linguistically Gifted Amateur}

Roger Williams was a theologian with no formal linguistic training, but his
orthography is the most phonetically sophisticated of the 17th-century
sources. His unique contribution is the systematic use of diacritics for
stress: acute (á) for primary stress, grave (à) for low/falling tone, and
circumflex (â) for rising-falling tone. This is a major advantage over
Eliot's Massachusett Bible (1663), which uses no accent marks, making stress
reconstruction for Massachusett far more difficult. However, Williams'
English phonological filter still produced systematic gaps: he almost never
wrote preaspirated /h/, and his five-vowel system (a, e, i, o, u) must
represent 8--10 Narragansett vowel phonemes.

\subsubsection{Eliot (1666) --- The Systematic Missionary}

John Eliot developed the first standardized writing system for an Algonquian
language, working with Massachusett collaborators including Job Nesutan and
John Sassamon. His orthography was phonemically oriented for its time --- he
attempted to represent the sounds of Massachusett systematically, not just
through English ear-approximations. His key innovation was distinguishing
aspirated from unaspirated stops using English voiced/voiceless pairs (b/p,
d/t, k/g). However, this system was still filtered through English
phonological categories: Eliot's \textit{b, d, g} represent unaspirated lenis
stops (not voiced), and his \textit{p, t, k} represent aspirated fortis
stops. The distinction is phonetically real in Massachusett but mapped onto
an English voicing contrast that does not align perfectly.

\subsubsection{Edwards (1787) --- The Fluent L2 Speaker}

Jonathan Edwards Jr.\ is unique among the colonial recorders in being a
fluent L2 speaker of Mahican, having learned the language as a child in
Stockbridge, Massachusetts, where Mahican families were the majority
population. His training in Hebrew and Greek gave him explicit phonological
categories, and his childhood fluency gave him native-like perception. The
result is the most reliable colonial transcription in the entire SNEA corpus.
Edwards consistently writes /h/ where PA reconstruction predicts it, and his
vowel transcriptions are more stable than any other colonial source. His
\textit{Observations on the Language of the Muhhekaneew Indians} (1788)
includes comparative glossaries with Shawnee and Ojibwa, demonstrating a
comparative method that was far ahead of its time.

\subsubsection{Prince-Speck (1904) --- The Two-Stage Chain}

J.\ Dyneley Prince and Frank Speck represent a two-stage recording chain
(fieldworker $\rightarrow$ desk transcriber) that introduces compounding
transcription errors. Speck, a Boasian anthropologist with minimal linguistic
training, collected field notes from Fidelia Fielding, the last fluent
speaker of Mohegan-Pequot. Prince, a Neogrammarian linguist who never heard
the language spoken, converted Speck's notes into the published orthography.
Prince's Germanic training introduced specific biases: \textit{tsh} for /tʃ/
(following German \textit{tsch}), \textit{z} for /ts/ (German \textit{z}),
and \textit{Dehnungs-h} (lengthening h) for vowel length. The glossary also
preserves some of Fielding's own English-based spellings for ``sentimental
reasons'' (Prince \& Speck 1904:19), creating a three-way orthographic
competition between English colonial conventions, Fielding's vernacular
spelling, and Prince's German linguistic conventions.

\subsection{The Kopris (2014) Framework}

Craig Kopris' paper ``Les sons wyandots perçus par des oreilles étrangères''
(Wyandot sounds perceived by foreign ears, 2014) provides the exact
theoretical framework for the colonial recorder problem. Kopris demonstrates
that European recorders systematically misperceived and mis-transcribed
indigenous phonemes because their native phonological systems lacked the
necessary categories. His key insight is that the recorder's L1 is not a
source of random error but a systematic filter: the errors are predictable
given the recorder's native phoneme inventory.

Applied to the SNEA corpus, the Kopris framework predicts:

\begin{enumerate}
  \item \textbf{English recorders will merge Algonquian stop series.} English
    has a two-way laryngeal contrast (voiced vs.\ voiceless aspirated);
    Algonquian languages have a different contrast (plain vs.\ preaspirated,
    or fortis vs.\ lenis). The mapping is systematically distorted.

  \item \textbf{Vowel quality will be underdifferentiated.} English has more
    vowel qualities than the five-letter Latin alphabet can represent, but
    colonial recorders still used only five vowel letters. The mapping from
    Algonquian vowels to English letters is systematically ambiguous.

  \item \textbf{Stress and length will be confused.} English stress is
    dynamic (loudness-based); Algonquian stress may be tonal or length-based.
    Recorders who mark stress (Williams) may be indirectly marking length,
    and vice versa.

  \item \textbf{Preaspiration will be invisible.} English has no phonemic
    preaspiration, so English-speaking recorders will systematically fail to
    hear it --- unless they have training (Prince) or L2 fluency (Edwards)
    that provides the necessary category.
\end{enumerate}

The Kopris framework transforms the colonial recorder problem from a
liability into an asset: once we know the recorder's L1 system, we can
predict the systematic distortions and reverse them. This is the theoretical
foundation for the cross-source calibration methodology described in
Section~\ref{sec:calibration}.
